\subsection{Paper 1: Efficient Classical Simulation of Heuristic Peaked Quantum Circuits}

Link: https://arxiv.org/abs/2404.14493

\begin{enumerate}
    \item The paper studies \emph{peaked quantum circuits}, whose output probability distribution is strongly concentrated on one hidden bitstring, making the result easy to verify once the peak is known.
    
    \item The specific circuits considered are generated by first training a shallow circuit to produce a peak and then inserting a deep obfuscated mirror block of the form
    \[
        U U^\dagger \approx I,
    \]
    which should ideally cancel but is hidden by permutations, masking, angle transformations, and swap insertions.
    
    \item Although this obfuscation makes direct symbolic cancellation difficult, the circuit still contains an exploitable mirror structure, because the left and right halves remain approximate inverses of each other.
    
    \item The authors use tensor-network methods, representing the central part of the circuit as a Matrix Product Operator (MPO), whose computational cost is controlled by its bond dimension.
    
    \item The bond dimension measures how much information or entanglement must be carried across a cut in the one-dimensional tensor network, but it can also grow artificially when qubit lines are badly ordered by hidden permutations.
    
    \item The circuit is split at its temporal midpoint as
    \[
        C = C_L \, M \, C_R,
    \]
    with \(M\) initialized as the identity MPO, and gates from \(C_L\) and \(C_R\) are iteratively absorbed into \(M\).
    
    \item For an unobfuscated mirror circuit, this middle-out contraction keeps the MPO small because inverse gates cancel, but in the obfuscated case hidden swaps cause the MPO bond dimension to grow.
    
    \item The key innovation is an ``unswapping'' procedure, which greedily tests nearest-neighbor swaps on the bonds with largest bond dimension and accepts swaps that reduce the MPO size, effectively decomposing
    \[
        M = P_L \, \widetilde{M} \, P_R,
    \]
    where \(P_L\) and \(P_R\) are extracted permutations and \(\widetilde{M}\) is a lower-bond-dimension MPO.
    
    \item After each unswapping step, the extracted permutations are pushed into the remaining left and right circuits by rewiring and re-transpiling, so that subsequent gate absorption continues in the correct qubit ordering.
    
    \item Once the full circuit has been contracted, the resulting MPO is applied to
    \[
        |0\rangle^{\otimes N}
    \]
    to obtain a Matrix Product State (MPS), from which samples are drawn efficiently; in the largest reported case, this recovers the hidden peak of a 56-qubit circuit with 1917 two-qubit gates in about one hour on a single GPU.
\end{enumerate}

\subsection{Paper 2: Heuristic Quantum Advantage with Peaked Circuits}

Link: https://arxiv.org/abs/2510.25838

